\subsection{P-17. Balanced two-pair rank bounds from coefficient cuts and $\mathrm{SL}_2$ symmetry}
\label{P-17}

\paragraph{Statement and scope.}
Fix distinct nodes $a_1,\ldots,a_n$ in a field $E$.
Let $\mathcal I_2$ be the vanishing ideal of evaluations of rational
functions with numerator and denominator degrees at most two.
All assertions below hold in every characteristic, with field
extensions used as proof devices when necessary.

Every nonzero multilinear polynomial in $\mathcal I_2$ on at most
eight variables has ROABP width at least four in every variable
order. The same conclusion holds for every nonzero homogeneous
multilinear cubic in $\mathcal I_2$ on at most ten variables,
and in any number of variables for the decomposable cubic class
specified below. In any number of variables, a nonzero multilinear
width-three polynomial whose multihomogenization is
$\mathrm{SL}_2$-invariant is also hit by balanced two-pair
evaluation, indeed by I-37's ordinary centered two-pair map.

These are statements about the balanced rational generator.
They do not identify it with the ordinary unweighted centered SV
map outside the separately stated invariant subclass. The
proof here leaves a common-translation-invariant nine-variable
quartic. The subsequent I-40 closes that residual, and I-41
proves the full multilinear conclusion through ten variables.
The unrestricted case beyond ten variables remains open.

\paragraph{Known ideal input and the decomposable class.}
Put
\[
 M=\begin{pmatrix}1&\cdots&1\\a_1&\cdots&a_n\\a_1^2&\cdots&a_n^2\end{pmatrix},
 \qquad K=\ker M\subseteq E^n.
 \tag{P17.1}
\]
For a three-element set $S$, use a consistent ordered Vandermonde
factor $\Delta_S=\det M_S$. Its sign convention is immaterial as
long as it is used consistently. Hu, van Melkebeek, and Morgan
(2024), Definition~9.5, p.~55, and Theorem~9.13, pp.~61--62
(P-B20), give the following exact input. A homogeneous multilinear
cubic
\[
 f=\sum_{|S|=3}c_S X_S
\]
lies in $\mathcal I_2$ if and only if its decoded alternating tensor
\[
 W=\sum_{|S|=3}\frac{c_S}{\Delta_S}e_S
       \quad\hbox{belongs to }\quad\bigwedge^3 K.
 \tag{P17.2}
\]
The same theorem says that the multilinear homogeneous degrees in
$\mathcal I_2$ lie in the interval $3\le d\le n-3$. We use these
published characterizations, rather than reprove them.

The decomposable class consists of the nonzero tensors
$W=p\wedge q\wedge v$ in (P17.2). If $Z$ is the $3\times n$
matrix with these three rows, the corresponding polynomial is,
by Cauchy--Binet,
\[
 f=\det\bigl(M\operatorname{diag}(X_1,\ldots,X_n)Z^{\mathsf T}\bigr),
 \qquad MZ^{\mathsf T}=0.
 \tag{P17.3}
\]
Every three columns of $M$ are independent. In particular, every
nonzero vector of $K$ has at least four nonzero coordinates.

\paragraph{1. A three-block cut obstruction.}
Let $f\in\mathcal I_2$ be a nonzero homogeneous multilinear cubic.
For any variable partition $U\sqcup V$, its coefficient matrix is
block diagonal after grouping row and column monomials by degree.
The block $C_d$ has left degree $d$ and right degree $3-d$.
We claim that three consecutive nonzero blocks $C_0,C_1,C_2$, or
$C_1,C_2,C_3$, force total coefficient rank at least four.

Suppose instead that $C_0,C_1,C_2$ are nonzero and the total rank
is at most three. Then $C_3=0$ and each of those three blocks has
rank one. Consequently
\[
 f=A(V)+L(U)B(V)+Q(U)N(V),
 \tag{P17.4}
\]
where $L,N$ are nonzero linear forms, $B,Q$ are nonzero homogeneous
quadratics, and $Q$ is multilinear. The pole-extraction identity
proved in I-37 gives
$\partial_i f\in\mathcal I_1$, where $\mathcal I_1$ denotes the
balanced one-pair ideal. The ideal $\mathcal I_1$ is prime and
contains no nonzero linear form: it is the kernel of a rational
parameter substitution, and the distinct simple poles
$1/(a_i-y)$ are linearly independent. Thus it contains no
nonzero product of two linear forms.

For every constant direction $b$ on $U$ with $b(L)=0$,
\[
 b(f)=b(Q)N\in\mathcal I_1,
\]
so $b(Q)=0$. Write $L=\ell^{\mathsf T}X_U$ and let $T$ be the
symmetric zero-diagonal coefficient matrix of the multilinear
quadratic $Q$, so that $b(Q)=(Tb)^{\mathsf T}X_U$ in every
characteristic. We have $\ker\ell^{\mathsf T}\subseteq\ker T$.
Therefore $T=u\ell^{\mathsf T}$ for a column $u$; symmetry implies
$u=c\ell$. Its zero diagonal then implies $c\ell_i^2=0$ for every
$i$, hence $c=0$, since $\ell\ne0$. Thus $T=0$ and $Q=0$, a
contradiction. This avoids division by two. Exchanging $U,V$
proves the claim for $C_1,C_2,C_3$.

\paragraph{2. Every decomposable cubic has a cut of rank at least four.}
Take (P17.3) and fix any variable order. Let $z_i\in E^3$ denote
column $i$ of $Z$. Choose the first prefix $U$ whose columns span
a two-dimensional space. Immediately before its last column,
the prefix span has dimension one. The remaining suffix $V$
has column rank at least two. Otherwise a nonzero linear
functional on $E^3$ annihilating the previous prefix and the
suffix would yield a nonzero row combination of $Z$ supported
only on that last prefix coordinate. This contradicts the
minimum support four of nonzero vectors in $K$.

If the suffix has rank three, there are nonzero monomial
coefficients of left degrees zero, one, and two. Indeed, choose
respectively three independent suffix columns, one nonzero
prefix column and two completing suffix columns, or two
independent prefix columns and one completing suffix column.
The Vandermonde factors in (P17.2) do not vanish. Part~1 therefore
forces coefficient rank at least four.

It remains to consider the case in which both projections of
the row space of $Z$ have rank two. Its kernels under these two
projections are distinct lines. Choose a row basis $p,q,v$ in
which $p$ is supported on $U$ and $q$ on $V$. Both $p$ and $q$
belong to $K$ and have at least four nonzero coordinates.
Moreover $v_U$ is not proportional to $p_U$, and $v_V$ is not
proportional to $q_V$, since the two projection ranks are two.

For $i\in U$ and $j,k\in V$, the left-degree-one coefficients
are, up to one consistent sign,
\[
 p_i(q_jv_k-q_kv_j)(a_j-a_k)(a_i-a_j)(a_i-a_k).
 \tag{P17.5}
\]
Evaluation of a polynomial of degree at most two at the nonzero
coordinates of $p$, followed by multiplication by those $p_i$,
is injective, because their nodes include at least four distinct
points. It follows that
\[
 \operatorname{rank}C_1
 =\dim\operatorname{span}\{
      (T-a_j)(T-a_k):q_jv_k-q_kv_j\ne0\}.
 \tag{P17.6}
\]
There are at least two distinct pairs with nonzero coefficients
in this set. If only one pair were nonzero, the decomposable
nonzero two-vector $q_V\wedge v_V$ would be supported on that
pair, so its two-dimensional row space, including $q_V$, would
be supported on only two coordinates. This contradicts the
support bound for $q$. Two distinct monic quadratics with
these distinct two-element root sets are linearly independent.
Thus $\operatorname{rank}C_1\ge2$.

The symmetric argument, using the nonzero coordinates of $q$
and the independent pair $p_U,v_U$, gives
$\operatorname{rank}C_2\ge2$. The degree blocks have disjoint
row and column indices, so the total coefficient rank is at
least four. This proves the decomposable assertion in every
order and for every $n$.

\paragraph{3. At most seven variables.}
For $n\le5$, the degree interval supplied by Theorem~9.13 is
empty. For $n=6$, the only possible degree is three and
$\dim K=3$, so every nonzero three-vector is decomposable.
For $n=7$, the possible degrees are three and four and
$\dim K=4$; every nonzero three-vector in a four-dimensional
space is decomposable. Thus part~2 covers every cubic.

Multilinear reversal
\[
 f^{\mathrm{rev}}(X)=\Bigl(\prod_iX_i\Bigr)f(X_1^{-1},\ldots,X_n^{-1})
 \tag{P17.7}
\]
preserves the balanced ideal, since reciprocal rational functions
still have numerator and denominator degrees at most two.
It preserves each coefficient-cut rank by complementing the
row and column subset indices, as already used in P-16. For
$n=7$, it turns a homogeneous quartic into a cubic, covering
that case as well.

Finally the ideal decomposes by homogeneous degree. The lowest
nonzero homogeneous component of a polynomial has ROABP width
no greater than the original in the same order (Hu et al., 2024,
Proposition~8.3, p.~42). Hence a nonhomogeneous multilinear
annihilator of width at most three would have a nonzero cubic
or quartic component of width at most three, contradicting the
preceding cases. This proves the full seven-variable assertion.

\paragraph{4. The additional eight-variable cubic case.}
Let $n=8$ and suppose that a cubic annihilator has width at most
three in a fixed order. Take the first three variables as $U$
and the remaining five as $V$. The restriction map
$K\to E^V$ is an isomorphism: its kernel would be a vector of
$K$ supported on three coordinates, and both spaces have
dimension five. Its third exterior power is also an
isomorphism. Therefore $C_0\ne0$ for every nonzero $W$.

We first show that $C_2=0$. If $C_1=0$, fix $i\in U,k\in V$.
The quadratic $\partial_i f\in\mathcal I_1$ has both zero row
sum and zero simple-pole residue at node $k$. Since $C_1=0$,
these two equations involve only the two possible indices
$j\in U\setminus\{i\}$:
\[
 \sum_j c_{ijk}=0,\qquad
 \sum_j\frac{c_{ijk}}{a_k-a_j}=0.
 \tag{P17.8}
\]
Their two distinct reciprocal weights make the coefficient
matrix invertible. Thus all these $c_{ijk}$ vanish and $C_2=0$.
If instead $C_1\ne0$ and $C_2\ne0$, part~1 contradicts the
rank bound because $C_0\ne0$. This proves $C_2=0$ in all cases.
The zero row-sum identity for a left pair then gives $C_3=0$,
since the sole possible third left index is its only remaining
term. The row-sum and residue identities used here follow
by substituting $t+z/(a_i-y)$ into the relevant quadratic
and taking the coefficient of $tz$ and the simple-pole
residues of the coefficient of $z^2$; no division by two is
used.

Put $K_0=K\cap E^V$, which has dimension two, and choose a
complement $J$ in $K$. The projection $J\to E^U$ is an
isomorphism. In the decomposition of $\bigwedge^3(K_0\oplus J)$,
the highest left-coordinate degree of $\bigwedge^3J$ is three,
and that of $K_0\wedge\bigwedge^2J$ is two; these leading
projections are injective. The absence of all coefficients
with at least two left variables therefore forces
\[
 W\in\bigwedge^2K_0\wedge J.
 \tag{P17.9}
\]
Because $\dim K_0=2$, every nonzero tensor in (P17.9) is a
wedge of a basis of $K_0$ with one additional vector. Thus
$W$ is decomposable. Part~2 now contradicts the width bound.
This proves the eight-variable cubic assertion.

\paragraph{5. The nine- and ten-variable cubics from two end anchors.}
The argument at the start of part~4 holds for any number of
variables. If $U$ is any three-element set and $V$ its complement,
the projection $K\to E^V$ is an isomorphism: it is injective by
the three-column independence of $M$, and both spaces have
dimension $n-3$. Thus the block $C_0$ is nonzero for every
nonzero cubic tensor $W$.

If this cut has total coefficient rank at most three, the same
two-equation argument (P17.8) and part~1 force $C_2=C_3=0$.
For $K_U=\ker(K\to E^U)$, choose a complement projecting
isomorphically to $E^U$. The leading-projection argument of
part~4 then gives
\[
 W\in\bigwedge^2K_U\wedge K.
 \tag{P17.12}
\]
This includes its possible component in $\bigwedge^3K_U$;
no decomposability is claimed for arbitrary $n$.

Now let $n=9$ and fix a putative width-three order. Apply this
anchor to its first three and its last three variables. The
corresponding kernels $K_L,K_R$ each have dimension three.
Their intersection is zero, since a vector in it is supported
on the middle three variables and belongs to $K$. Therefore
\[
 K=K_L\oplus K_R.
\]
The first anchor requires at least two factors from $K_L$ in
the third exterior power, and the second requires at least
two from $K_R$. These are disjoint summands in
\[
 \bigwedge^3(K_L\oplus K_R)
 =\bigoplus_{j=0}^3
      \bigwedge^jK_L\otimes\bigwedge^{3-j}K_R.
\]
Their intersection is zero, contradicting $W\ne0$.
This proves the nine-variable cubic assertion in every order.

For $n=10$, the two end kernels have dimension four and their
intersection $P$ is the one-dimensional kernel supported on the
middle four coordinates. Write
\[
 K=P\oplus A\oplus B,\qquad
 K_R=P\oplus A,\quad K_L=P\oplus B,
 \tag{P17.13}
\]
where $A$ and $B$ each have dimension three. The first-coordinate
projection maps $A$ isomorphically to the first three coordinates,
and the last-coordinate projection does the same for $B$ and the
last three. In the exterior-power decomposition associated with
(P17.13), the first anchor requires at least two factors from
$P\oplus B$, and the last anchor at least two from $P\oplus A$.
Since $\dim P=1$, their intersection is exactly
$P\wedge A\wedge B$. Choose $0\ne p\in P$ and bases
$A_1,A_2,A_3$, $B_1,B_2,B_3$ whose respective outer projections
are the coordinate bases. Then
\[
 W=p\wedge\sum_{i,j=1}^3 Z_{ij}A_i\wedge B_j,
 \qquad Z\ne0.
 \tag{P17.14}
\]
Every one of the four middle coordinates of $p$ is nonzero,
by the minimum support four of vectors in $K$.

Take the cut after the first four variables. Fix $Z_{ij}\ne0$.
For each middle coordinate $k$, the coefficient of the triple
consisting of first coordinate $i$, middle coordinate $k$, and
last coordinate $j$ is, up to its ordered sign,
$p_kZ_{ij}\Delta_{\{i,k,j\}}\ne0$. Using the first and second
middle coordinates gives respectively a nonzero entry in $C_2$
and in $C_1$. If $C_0\ne0$, part~1 contradicts the cut-rank
bound. If $C_0=0$, the projection from $K$ to the last six
coordinates has a one-dimensional kernel $Eq$, consisting of
the vectors supported on the first four coordinates. Its
induced map on $K/Eq$ is injective, so the kernel of its third
exterior power is exactly $q\wedge\bigwedge^2K$. Consequently
$W$ is divisible by both $p$ and $q$. The two lines are distinct,
because their supports intersect in only one coordinate and
each nonzero vector of $K$ has support at least four. Thus
\[
 W\in\left(p\wedge\bigwedge^2K\right)
       \cap\left(q\wedge\bigwedge^2K\right)
   =p\wedge q\wedge K.
\]
The nonzero tensor $W$ is therefore decomposable, contradicting
part~2. This proves the ten-variable cubic assertion in every
order and characteristic.

\paragraph{6. An invariant multilinear width-three tensor has a singular affine representation.}
Work over an algebraic closure; polynomial identities and
coefficient-cut ranks are unaffected by this field extension.
We use the characteristic-free weight decomposition for the split
torus $\mathbb G_m$: its rational representations are
$\mathbb Z$-graded vector spaces. The exact primary reference is
SGA~3, Expos\'e~I, Proposition~4.7.3 (G-B02), whose proof is the
orthogonal-projector decomposition of its comodule map. This
established torus fact is used here, not reproved.

For a multilinear polynomial $f$ in an even number $n$ of variables,
write its multihomogenization as
\[
 \widetilde f((U_i,V_i)_{i=1}^n)
   =\Bigl(\prod_i U_i\Bigr)f(V_1/U_1,\ldots,V_n/U_n).
 \tag{P17.10}
\]
Suppose that this tensor is invariant under the simultaneous
$\mathrm{SL}_2$ action on all pairs, and that its coefficient-cut
ranks in some order are at most three. For a cut after $i$
variables, its left Schmidt space is the image of contraction
with the dual of the right tensor factors. Invariance makes
this image an $\mathrm{SL}_2$ submodule of $(E^2)^{\otimes i}$,
and its dimension is exactly the coefficient-cut rank $r_i$.
Restrict this Schmidt space to the diagonal torus
$\operatorname{diag}(t,t^{-1})$. The ambient tensor power has
integer character weights $i,i-2,\ldots,-i$. For odd $i$, all
these weights are odd and none is zero. The Weyl element
\[
 w=\begin{pmatrix}0&-1\\1&0\end{pmatrix}\in\mathrm{SL}_2
\]
conjugates the torus parameter $t$ to $t^{-1}$, and hence gives
an isomorphism between the weight-$r$ and weight-$-r$ subspaces.
The nonzero integer weights pair without fixed points, so the
Schmidt-space dimension is even. These are torus characters,
not eigenvalues reduced modulo the field characteristic; the
argument also holds in characteristic two. Since $r_i\le3$
and $f\ne0$, we obtain $r_i=2$ at every odd cut.

Use the standard minimal tensor-train factorization with bond
sizes $r_i$ (P-B17, Theorem~2.1), with $r_0=r_n=1$.
Each local core is affine in its one variable and has matrix
shape $r_{i-1}\times r_i$. One of its two adjacent cuts is odd,
so one side has dimension at most two. Padding these cores
consistently to $3\times3$ matrices gives an affine width-three
word whose every variable-dependent factor has rank at most
two. I-37 therefore hits $f$ with two centered pairs, both in
its ordinary and its rational formulations. This proves the
stated invariant subclass for every even $n$.

\paragraph{7. The eight-variable quartic in every characteristic.}
Suppose that a homogeneous multilinear quartic $f\in\mathcal I_2$
on eight variables has width at most three. Balanced rational
evaluation is closed under adding a constant, so every coefficient
of $f(X+t\mathbf1)$, as a polynomial in $t$, lies in
$\mathcal I_2$. Its coefficient of $t^j$ has homogeneous degree
$4-j$. The known minimum degree three therefore gives
\[
 f(X+t\mathbf1)=f(X)+t g(X),\qquad
 g\in\mathcal I_2\text{ homogeneous of degree three}.
 \tag{P17.11}
\]
If $g\ne0$, it has width at most three in the same order.
Indeed translating each univariate factor preserves the cut
ranks over $E(t)$; multiply the coefficient matrices by $t^{-1}$
and specialize $t^{-1}$ to zero. All minors of size four still
vanish, so the limiting coefficient matrices of $g$ have rank
at most three. Part~4 excludes such a nonzero $g$. Hence
$f(X+t\mathbf1)=f(X)$.

The reversal $f^{\mathrm{rev}}$ is also a homogeneous quartic
annihilator of width at most three. Applying the same argument
to it proves its common-translation invariance. On the
multihomogenization (P17.10), common translation is one
unipotent subgroup of $\mathrm{SL}_2$, while reversal conjugates
it to the opposite unipotent subgroup. Those two subgroups
generate $\mathrm{SL}_2$. Thus $\widetilde f$ is invariant,
and part~6 contradicts its membership in $\mathcal I_2$.

The only possible homogeneous degrees on eight variables are
three, four, and five. Part~4 excludes degree three, reversal
excludes degree five, and the preceding argument excludes
degree four. Passing to a lowest nonzero homogeneous component
as in part~3 proves the full eight-variable conclusion.\hfill$\square$

\paragraph{The remaining balanced-ideal gate.}
The degree-zero-block anchor holds uniformly for any $n$.
The extra four-coordinate cut resolves the one-dimensional
intersection of the two end kernels at $n=10$. At $n=11$ their
intersection has dimension two, so that argument no longer
forces the same decomposable tensor. The unrestricted cubic
case from eleven variables onward remains open here.

On nine variables, the possible homogeneous degrees are three
through six. Part~5 and reversal exclude degrees three and six.
A putative width-three annihilator therefore reduces, after
homogeneous-component extraction and possibly reversal, to a
quartic. The translation-leading argument of part~7, now using
the nine-variable cubic bound, makes that quartic invariant under
common translation. Its multihomogenization has nonzero torus
weight, so the invariant-tensor argument does not automatically
apply. This common-translation-invariant nine-variable quartic
is the smallest case left by the proof.

The nine-variable residual of this proof is subsequently closed
by I-40, and I-41 proves the full multilinear balanced two-pair
theorem through ten variables. The unrestricted case beyond ten
variables remains a separate gate. The proof does not infer
cut-rank bounds from a fixed linear change of coordinates to a
Hankel ideal; such a change need not preserve ROABP coefficient
cuts. The all-characteristic odd-weight argument and the
ten-variable cubic extension were supplied independently by
the isolation agent during the proof review,
replacing the author's initial characteristic-zero
complete-reducibility argument. No historical priority claim
is made.
